\section{System Model}
\label{sec:system}

We consider a grant-free uplink from $K$ single-antenna terminals to a
LEO satellite with a uniform planar array (UPA) of $M=M_yM_z$ elements.
Each access block contains $Q$ AFDM pilot frames with $N$ useful samples
per frame, over which activities, channel geometry, and path coefficients
are held fixed. With common block timing and compensation of a reference
delay $\tau_c$, differential delays and Doppler shifts are handled at
the satellite without user-specific pre-compensation \cite{shen2025}.
The receiver estimates activities, channels, inter-frame phases, and
noise power from the superimposed pilots.

Let $\alpha_k\in\{0,1\}$ denote terminal activity. Independent
activations with probability $p$ are conditioned on at least one
transmission, giving the marginal $\rho=p/[1-(1-p)^K]$ supplied to
the receiver.

\subsection{Terrestrial-Satellite Link Model}

The propagation geometry is specified by the satellite position
$\mathbf{s}$, its velocity $\mathbf{v}_{\rm sat}$, and the position
$\mathbf{g}_k$ of terminal $k$.
With carrier frequency $f_c$ and propagation speed $c_0$, define
\begin{equation}
 \mathbf{u}_k=\frac{\mathbf{g}_k-\mathbf{s}}
 {\norm{\mathbf{g}_k-\mathbf{s}}},
 \qquad
 \nu_k=-\frac{f_c}{c_0}\mathbf{v}_{\rm sat}^{\mathsf T}\mathbf{u}_k.
 \label{eq:geometry}
\end{equation}
With the convention $s_{\rm RF}(t)=\Re\{x(t)e^{-j2\pi f_ct}\}$,
a path contributes $x(t-\tau(t))e^{j2\pi f_c\tau(t)}$ up to its gain,
giving the Doppler factor $e^{j2\pi\nu_kt}$ in \eqref{eq:geometry}.
We retain this carrier phase variation while holding envelope delays
and array directions fixed over the block.

After compensation of the common reference delay $\tau_c$, the residual
delay of path $\ell$ is given by
\begin{equation}
 \tau_{k\ell}=\frac{\norm{\mathbf{g}_k-\mathbf{s}}}{c_0}-\tau_c
             +\delta_\ell\tau_{{\rm spread},k},
 \label{eq:path-delay}
\end{equation}
where $\tau_c$ is the propagation time to the closest ground point in the
beam footprint, and $\delta_\ell\tau_{{\rm spread},k}$ is the excess
multipath delay. We adopt a common spatial response and a common Doppler
shift for the paths of each terminal
\cite{li2023,huang2026}, while the path gains and excess delays remain
path-dependent. This common Doppler allows the phase evolution across
frames to be separated from the sum of the path contributions.

Let $\mu_{y,k}=\mathbf e_y^{\mathsf T}\mathbf u_k$ and
$\mu_{z,k}=\mathbf e_z^{\mathsf T}\mathbf u_k$, where $\mathbf e_y$ and
$\mathbf e_z$ are the array axes. For element position
$\mathbf s-n_yd_y\mathbf e_y-n_zd_z\mathbf e_z$, the array response is
\begin{equation}
 \begin{gathered}
 [\mathbf{a}_k]_m=
 \exp\!\left[j2\pi\left(
 \frac{d_y}{\lambda}n_y\mu_{y,k}
 +\frac{d_z}{\lambda}n_z\mu_{z,k}\right)\right],\\
 m=n_z+M_zn_y,\qquad \norm{\mathbf{a}_k}^2=M,
 \end{gathered}
 \label{eq:upa}
\end{equation}
where $\lambda=c_0/f_c$, $0\leq n_y<M_y$, and $0\leq n_z<M_z$.
The unknown array response is included in the channel to be estimated.

\subsection{Transmission Scheme}

Each active terminal sends known pilots that vary across the $Q$
frames, providing measurements of a shared channel after the
inter-frame phase is separated from the single-frame operator.
Let $T_s=1/(N\Delta f)$ be the sampling interval, where $\Delta f$
denotes the subcarrier spacing. The DAFT matrix is \cite{bemani2022}
\begin{equation}
 \begin{gathered}
 [\mathbf{F}_{\rm A}]_{rt}=\frac{1}{\sqrt N}
 e^{-j2\pi(c_2r^2+rt/N+c_1t^2)},\\
 r,t=0,\ldots,N-1.
 \end{gathered}
 \label{eq:daft}
\end{equation}
For pilot vector $\mathbf p_{qk}$ with $[\mathbf p_{qk}]_n=P_{qkn}$,
the useful time samples are
$\mathbf s_{qk}=\mathbf F_{\rm A}^{\mathsf H}\mathbf p_{qk}$.
A chirp-periodic prefix (CPP) of length $N_{\rm CPP}$ extends these samples as
\begin{equation}
 \begin{gathered}
 s_{qk}[t]=e^{j2\pi c_1(t^2-\langle t\rangle_N^2)}
 s_{qk}[\langle t\rangle_N],\\
 -N_{\rm CPP}\leq t<0,
 \end{gathered}
 \label{eq:cpp}
\end{equation}
where $\langle x\rangle_N=x-N\lfloor x/N\rfloor$.
We choose $N_{\rm CPP}T_s\geq\max_{k,\ell}\tau_{k\ell}$ and
$2Nc_1\in\mathbb N_+$. The frame duration is
$T_{\rm frame}=(N+N_{\rm CPP})T_s$.
The prefix covers the residual propagation delays, including delays
longer than the useful interval $NT_s$, so that each useful window
contains samples from the corresponding transmitted frame.

Fractional delay is retained in constructing the AFDM effective channel
\cite{huang2026}. Let $w_\xi=\langle\xi\rangle_N$ and
$\zeta_n(w)=\lfloor2c_1w+n/N\rfloor$. We write the AFDM basis as
\begin{equation}
 \varphi_n(\xi)=\frac{1}{\sqrt N}
 e^{j2\pi[c_2n^2+c_1\xi^2+nw_\xi/N-\zeta_n(w_\xi)w_\xi]}.
 \label{eq:continuous-waveform}
\end{equation}
It agrees with the inverse DAFT at integer useful-sample times and
extends across prefix periods according to \eqref{eq:cpp}.
For delay $d=\tau/T_s$, the delayed basis
$[\mathbf B_d]_{tn}=\varphi_n(t-d)$ gives the DAFT-domain path
operator at useful-window start index $g$ as
\begin{equation}
 \begin{aligned}
 \mathbf G(d,\nu;g)&=\mathbf F_{\rm A}\mathbf D_{\nu,g}\mathbf B_d,\\
 \mathbf D_{\nu,g}&=\diag\{e^{j2\pi\nu T_s(g+t)}\}_{t=0}^{N-1}.
 \end{aligned}
 \label{eq:path-operator}
\end{equation}
With $\zeta_{tn}=\zeta_n(\langle t-d\rangle_N)$, expanding this
operator gives
\begin{equation}
 \begin{split}
 G_{rn}
 &=\frac{e^{j2\pi[c_1d^2+c_2(n^2-r^2)-nd/N+\nu T_sg]}}{N}\\
 &\quad\times\sum_{t=0}^{N-1}
 e^{j2\pi[(n-r)/N-2c_1d+\nu T_s]t}
 e^{j2\pi\zeta_{tn}d}.
 \end{split}
 \label{eq:path-entry-expansion}
\end{equation}
The last factor equals one for integer $d$, while continuous $d$ and
$\nu$ retain fractional delay and Doppler effects in the channel
coefficients without assigning paths to delay--Doppler grid points.

For integer $d$, \eqref{eq:path-operator} reduces to
\begin{align}
 \mathbf G(d,\nu;g)
 &=\mathbf F_{\rm A}\mathbf D_{\nu,g}\bm\Omega_d
   \bm\Pi_{d\bmod N}\mathbf F_{\rm A}^{\mathsf H},
 \label{eq:integer-operator}\\
 [\bm\Omega_d]_{tt}
 &=e^{j2\pi c_1[(t-d)^2-\langle t-d\rangle_N^2]},
 \label{eq:integer-prefix-phase}
\end{align}
where $[\bm\Pi_l\mathbf s]_t=s[\langle t-l\rangle_N]$.
The full delay enters the prefix phase, and its remainder determines
the cyclic shift.

At the reference frame, the channel of terminal $k$ is given by
\begin{equation}
 \begin{aligned}
 \mathbf G_k&=\sum_\ell g_{k\ell}
 \mathbf G(\tau_{k\ell}/T_s,\nu_k;N_{\rm CPP}),\\
 H_{kmrn}&=[\mathbf a_k]_m[\mathbf G_k]_{rn},
 \end{aligned}
 \label{eq:channel}
\end{equation}
where $g_{k\ell}$ is the path gain after large-scale power compensation
and $H_k\in\mathbb C^{M\times N\times N}$ indexes antennas, DAFT
outputs, and DAFT inputs. The common spatial response makes its
$M\times N^2$ reshaping rank at most one, which is also used for
noise estimation in Appendix~\ref{app:frontend}.

For $Q\leq KN$, pilots are formed from selected rows and permuted columns
of a $KN$-point discrete Fourier transform (DFT). With distinct row indices
$r_q$, a column permutation $c_{kn}$, and independent uniform phases
$\phi_{kn}$, set
\begin{equation}
 \begin{gathered}
 P_{qkn}=\sqrt{\frac{E_p}{N}}
 e^{-j2\pi r_qc_{kn}/(KN)}e^{j\phi_{kn}},\\
 \sum_{n=0}^{N-1}|P_{qkn}|^2=E_p.
 \end{gathered}
 \label{eq:pilot}
\end{equation}
The common terminal Doppler introduces the phase increment
$\beta_k=\operatorname{wrap}(2\pi\nu_kT_{\rm frame})$ between consecutive
frames, where $\operatorname{wrap}$ maps an angle to $[-\pi,\pi)$.
To obtain the channel in frame $q$, its useful-window start in
\eqref{eq:path-operator} is set to
$g_q=N_{\rm CPP}+q(N+N_{\rm CPP})$. Since
$\mathbf D_{\nu_k,g_q}=e^{jq\beta_k}\mathbf D_{\nu_k,N_{\rm CPP}}$,
the channel operator satisfies
\begin{equation}
 \begin{aligned}
 \mathbf G_k^{(q)}
 &=\sum_\ell g_{k\ell}\mathbf G(\tau_{k\ell}/T_s,\nu_k;g_q)\\
 &=e^{jq\beta_k}\mathbf G_k.
 \end{aligned}
 \label{eq:twc-frame-channel}
\end{equation}
The prefix condition removes interference from adjacent frames, while
the blockwise constant gains and common Doppler preserve $\mathbf G_k$,
so the frame-dependent phase can be included in the pilot operator as
\begin{equation}
 A_{qkn}=P_{qkn}e^{jq\beta_k},\qquad q=0,\ldots,Q-1.
 \label{eq:effective-pilot}
\end{equation}
The delay and within-frame Doppler effects remain in $H_k$.

Let $X_k=\alpha_kH_k$ denote the channel multiplied by its activity
indicator, so that $X_k$ vanishes for an inactive terminal.
After prefix removal and DAFT demodulation, the received signal can be expressed as
\begin{equation}
 \begin{gathered}
 Y_{qmr}=\sum_{k=1}^K\sum_{n=0}^{N-1}A_{qkn}X_{kmrn}+W_{qmr},
 \\
 W_{qmr}\sim\CN(0,\sigma^2),
 \end{gathered}
 \label{eq:observation-tensor}
\end{equation}
with conditionally independent noise entries.
Using the indexing
\begin{equation}
 \begin{aligned}
 [\mathbf Y]_{q,mN+r}&=Y_{qmr},\\
 [\mathbf A]_{q,(k-1)N+n}&=A_{qkn},\\
 [\mathbf X]_{(k-1)N+n,mN+r}&=X_{kmrn},
 \end{aligned}
 \label{eq:mmv-indexing}
\end{equation}
the MMV model can be written as
\begin{equation}
 \begin{gathered}
 \mathbf Y=\mathbf A\mathbf X+\mathbf W,\\
 \mathbf Y\in\mathbb C^{Q\times MN},\qquad
 \mathbf A\in\mathbb C^{Q\times KN},\\
 \mathbf X\in\mathbb C^{KN\times MN}.
 \end{gathered}
 \label{eq:mmv}
\end{equation}
All $MN$ measurement vectors share the activity pattern, while
increasing $Q$ adds observation rows without changing $\mathbf X$.
This requires a longer block $T_{\rm B}=QT_{\rm frame}$, over which
delays, directions, gains, and Doppler shifts are approximated as
constant, with carrier phase evolution retained in
\eqref{eq:twc-frame-channel}.

To assess this block approximation, as in \cite{shen2025}, consider a
fixed ground terminal and a circular orbit of altitude $h_{\rm sat}$.
Differentiating \eqref{eq:geometry} gives
\[
 |\dot\nu_k|\leq B_\nu
 =\frac{f_c}{c_0}\left(\frac{v_{\rm sat}^2}{h_{\rm sat}}
       +\frac{v_{\rm sat}^2}{R_{\rm E}+h_{\rm sat}}\right),
\]
where $v_{\rm sat}=\|\mathbf v_{\rm sat}\|$ and $R_{\rm E}$ is the
Earth radius. With fixed $\tau_c$ and initial Doppler magnitude at
most $\nu_{\max}$, the delay change is bounded by
$(\nu_{\max}T_{\rm B}+B_\nu T_{\rm B}^2/2)/f_c$ and the phase error
from holding Doppler at its initial value by $\pi B_\nu T_{\rm B}^2$. For $N=16$,
$\Delta f=30$~kHz, $N_{\rm CPP}=336$, and $Q=16$, $T_{\rm B}=11.73$~ms.
Taking $R_{\rm E}=6371$~km, $f_c=2$~GHz, $h_{\rm sat}=600$~km,
$v_{\rm sat}=7.56$~km/s, and $\nu_{\max}=41$~kHz gives bounds of
$0.241\,\mu$s ($0.116T_s$), $8.10$~Hz, and $0.30$~rad for delay
and Doppler changes and the phase error, respectively. Thus additional frames
provide more observations at the cost of greater within-block
variation, with the phase-error bound growing quadratically.

For $Q<KN$, the linear system for each measurement vector is
underdetermined, so we exploit the shared activity pattern and the
distribution of the effective-channel coefficients to estimate $\mathbf X$.
The pilot matrix also affects the linear estimator: although the
original DFT pilots have orthogonal rows, the row inner products of
$\mathbf A$ contain the user-dependent factors
$e^{j(q-q')\beta_k}$, which generally remove row orthogonality.
The linear estimator therefore uses the full matrix
$\mathbf A\mathbf A^{\mathsf H}$.

From the observations, pilots $P$, and activity marginal $\rho$, the
complete receiver estimates $\widehat{\bm\beta}$ and $\widehat\sigma^2$
and constructs $\widehat{\mathbf A}=\mathbf A(\widehat{\bm\beta})$
as described in Appendix~\ref{app:frontend}.
For simulations, the signal-to-noise ratio (SNR)
is defined using the clean superposition $\mathbf Y_0=\mathbf A\mathbf X$:
\begin{equation}
 \sigma^2=\frac{\norm{\mathbf Y_0}_F^2}{QMN}
 10^{-\mathrm{SNR}_{\rm dB}/10}.
 \label{eq:snr}
\end{equation}
The parameters $(\mathbf A_{\rm c},\sigma_{\rm c}^2)$ held fixed during
message updates are $(\widehat{\mathbf A},\widehat\sigma^2)$ for the
complete receiver and the supplied $(\mathbf A,\sigma^2)$ for training.

\subsection{Problem Formulation}

Let $D=MN^2$ and let $q_{\rm tr}$ be the distribution of channel tensors
generated by \eqref{eq:channel} before activity masking, normalized by
the training-channel root mean square $s$.
We divide $(\mathbf Y,H_k,X_k,\mathbf W)$ by $s$ and
$(\sigma^2,\sigma_{\rm c}^2)$ by $s^2$, leaving the pilot operators
unchanged; the same symbols denote normalized quantities below, and
channel outputs are rescaled by $s$.

We retain the joint distribution of the $D$ coefficients in each user
block and approximate the prior across users and activity states by
the product of its marginals. Conditioned on $\mathbf A_{\rm c}$ and
$\sigma_{\rm c}^2$, this gives the posterior
\begin{equation}
 \begin{aligned}
 &q(d\mathbf X,\{dH_k\},\bm\alpha\mid
 \mathbf Y,\mathbf A_{\rm c},\sigma_{\rm c}^2)\\
 &\quad\propto
 e^{-\norm{\mathbf Y-\mathbf A_{\rm c}\mathbf X}_F^2/\sigma_{\rm c}^2}\\
 &\qquad\times\prod_{k=1}^{K}
 \delta_{\alpha_kH_k}(dX_k)q_{\rm tr}(dH_k)
 \rho^{\alpha_k}(1-\rho)^{1-\alpha_k},
 \end{aligned}
 \label{eq:working-posterior}
\end{equation}
where $\delta_{\alpha_kH_k}$ enforces $X_k=\alpha_kH_k$.
Under this posterior, the minimum mean-squared error (MMSE) target
and the activity decision are given by
\begin{equation}
 \begin{aligned}
 X_k^{\rm MMSE}&=\E_q[X_k\mid\mathbf Y,\mathbf A_{\rm c},\sigma_{\rm c}^2],\\
 p_k^{\rm Bayes}&=q(\alpha_k=1\mid\mathbf Y,\mathbf A_{\rm c},\sigma_{\rm c}^2),\\
 \alpha_k^{\rm Bayes}&=\mathbf1\{p_k^{\rm Bayes}\geq\tfrac12\}.
 \end{aligned}
 \label{eq:targets}
\end{equation}
The threshold assigns equal costs to missed detections and false
alarms. Since the likelihood couples users through their superimposed
pilots, marginalization requires integration over all user channels
and summation over their activities. We therefore use message passing
to separate linear estimation under the observation factor from
user-level estimation under the channel and activity factors.
